\scanpage{64}
% ll. 1–8
D'autre part, pour une situation projective donnée, il y a 4 choix de situations hyperpolygônales à lui donner naissance (suivant le choix d'un des deux pts fixe pour $\tau_0$, et pour $\tau_1$) on passe d'une telle situation \unsure{(numérique|(correspondante|hyperpolygonale} : $\lambda_0,\lambda_1,\ \beta=\zeta(\lambda_0-1)(\lambda_1-1)$ à une autre en remplaçant $\tau_0$ par $\frac{1}{\lambda_0}\tau_0$, ou $\tau_1$ par $\frac{1}{\lambda_1}\tau_1$, ou l'un et l'autre à la fois ---

% ll. 9–14
ce qui ne change pas $K_0,K_1,K'$ car $\tau_0,\tau_1$ et $\tau_0\tau_1$ sont multipliés chacun par un scalaire. \deleted{Donc en principe} Mais ces op. remplacent $(\lambda_0,\lambda_1)$ par $(\lambda_0^{-1},\lambda_1)$ resp $(\lambda_0,\lambda_1^{-1})$ resp $(\lambda_0^{-1},\lambda_1^{\unsure{-|-1}})$ qui \deleted{sont déjà différentes} donnent \inserted{``en général''} \unsure{quadruplement|4 systèmes|un système} de valeurs diff.,

% ll. 15–19
\deleted{tandis que les valeurs correspondantes de $\beta^{\unsure{=}}$, qui ne prend \deleted{au plus} \inserted{a priori} que 2 val. possible, d'après l'équ. donnée plus haut sur $\beta$}
--- ``en général'' signifiant que $\nu_0$ et $\nu_1$ sont $\ne2$, i.e. $\lambda_0,\lambda_1\ne-1$ (i.e. $\mu_0,\mu_1\ne-2$).

% ll. 20–22
Dans ce cas, il s'ensuit donc que pour \deleted{les (2 solutions éventuelles} \inserted{$\lambda_0,\lambda_1$ et $K$ fixés, \deleted{\unsure{quand}} $K\ne-2$} \unsure{NB.} l'équation en $\zeta$ a a priori 2 soln
% archive sheet number: 63
